\chapter{Setting Up Your Development Environment}
\label{app:env}

This appendix gets you from a fresh computer to a built, runnable \prg on Windows,
Linux, or macOS. The instructions follow the platform's own setup guides
\citep{prg32repo}. The common foundation on every operating system is the same:
install Espressif's ESP-IDF, the development framework for these \riscv
microcontrollers, then build the firmware for your chosen target. Pick your
operating system's section, then read the shared sections on QEMU and the build
commands.

\section{The one idea behind every setup}

Whatever your operating system, the toolchain is ESP-IDF version 5.3 or newer,
configured for the \code{esp32c3} and \code{esp32c6} targets---the first for the
QEMU emulator, the second for the physical board \citep{prg32repo,prg32qemu}. Once
ESP-IDF is installed and ``sourced'' (made active in your shell), the command
\code{idf.py} builds, flashes, and runs everything. Almost every setup problem
reduces to ESP-IDF not being active in the current shell.

\begin{prgwarn}
ESP-IDF must be activated in \emph{each new terminal} before \code{idf.py} works.
If you see \code{idf.py: command not found} or a missing \code{riscv32-esp-elf-gcc},
ESP-IDF is simply not sourced in that shell. This is the single most common
setup error, and the fix is always the same: re-run the export step
\citep{prg32repo}.
\end{prgwarn}

\section{macOS}

The platform's macOS path, tested on Apple Silicon, installs the prerequisites with
Homebrew and then ESP-IDF \citep{prg32repo}.

\begin{lstlisting}[style=bash,caption={macOS setup with Homebrew and ESP-IDF.}]
brew install git cmake ninja dfu-util ccache libusb python
cd $HOME
git clone -b v5.3 --recursive https://github.com/espressif/esp-idf.git
cd esp-idf
./install.sh esp32c3,esp32c6
. ./export.sh
\end{lstlisting}

A convenient habit is an alias so a single word activates ESP-IDF in any shell
\citep{prg32repo}.

\begin{lstlisting}[style=bash,caption={A quality-of-life alias for activating ESP-IDF.}]
alias get_idf=". $HOME/esp-idf/export.sh"
\end{lstlisting}

\section{Linux}

On Debian or Ubuntu, install the build prerequisites with the system package
manager, then ESP-IDF as above \citep{prg32repo}. Other distributions use their own
package names for the equivalent libraries.

\begin{lstlisting}[style=bash,caption={Linux (Debian/Ubuntu) prerequisites, then ESP-IDF.}]
sudo apt update
sudo apt install -y git wget flex bison gperf python3 python3-venv \
  python3-pip cmake ninja-build ccache libffi-dev libssl-dev \
  dfu-util libusb-1.0-0
cd $HOME
git clone -b v5.3 --recursive https://github.com/espressif/esp-idf.git
cd esp-idf
./install.sh esp32c3,esp32c6
. ./export.sh
\end{lstlisting}

To flash a real board, your user must be allowed to use the serial port. Add
yourself to the \code{dialout} group, then log out and back in
\citep{prg32repo}.

\begin{lstlisting}[style=bash,caption={Granting serial-port access on Linux.}]
sudo usermod -aG dialout "$USER"   # then log out and back in
\end{lstlisting}

\begin{prgaside}
ESP32\nobreakdash-C6 boards usually appear as \fname{/dev/ttyACM0}, while USB-bridge boards may
appear as \fname{/dev/ttyUSB0} \citep{prg32repo}. If flashing fails although the
board is visible, the cause is almost always the \code{dialout} membership not yet
taking effect---log out and back in after adding yourself to the group.
\end{prgaside}

\section{Windows}

The platform recommends a classroom path built on Visual Studio Code with the
Espressif ESP-IDF extension, which installs the toolchain for you, plus the
PlatformIO and Microsoft C/C++ and Python extensions \citep{prg32repo}. There is
also a standalone path using the official ESP-IDF Tools Installer, after which you
work from the ``ESP-IDF PowerShell'' shortcut it creates.

\begin{lstlisting}[style=bash,caption={Windows standalone path, run from the ESP-IDF PowerShell shortcut.}]
cd $HOME\Documents
git clone https://github.com/riscv-prg32/PRG32.git
cd PRG32
idf.py -B build-esp32c6 -D SDKCONFIG_DEFAULTS=sdkconfig.defaults set-target esp32c6
idf.py -B build-esp32c6 -D SDKCONFIG_DEFAULTS=sdkconfig.defaults build
\end{lstlisting}

\begin{prgwarn}
On Windows, use the ESP-IDF PowerShell or ESP-IDF Command Prompt, not a plain
terminal where \code{idf.py} has not been exported. If flashing fails, check Device
Manager for the board's serial port and pass it explicitly with \code{-p COMx}; if
the monitor cannot open the port, close every other program using it (Arduino
Serial Monitor, other ESP-IDF monitors, and so on) \citep{prg32repo}.
\end{prgwarn}

\section{Installing the QEMU emulator}

To run on the desktop, install Espressif's \riscv QEMU after ESP-IDF, then
re-activate ESP-IDF so the new tool is on the path \citep{prg32qemu}.

\begin{lstlisting}[style=bash,caption={Installing the RISC-V QEMU tool and re-activating ESP-IDF.}]
python $IDF_PATH/tools/idf_tools.py install qemu-riscv32
. $IDF_PATH/export.sh
\end{lstlisting}

Linux and macOS hosts also need a few native libraries for the QEMU display window
\citep{prg32qemu}.

\begin{lstlisting}[style=bash,caption={QEMU window libraries on Linux (top) and macOS (bottom).}]
sudo apt-get install -y libgcrypt20 libglib2.0-0 libpixman-1-0 libsdl2-2.0-0 libslirp0
brew install libgcrypt glib pixman sdl2 libslirp
\end{lstlisting}

\section{The build commands, both targets}

With ESP-IDF active, the same two pairs of commands appear throughout the book.
Keep the two build directories strictly separate, as warned in
Chapter~\ref{ch:platform} \citep{prg32qemu}.

\begin{lstlisting}[style=bash,caption={Building and flashing for the physical ESP32\nobreakdash-C6 board.}]
idf.py -B build-esp32c6 -D SDKCONFIG_DEFAULTS=sdkconfig.defaults set-target esp32c6
idf.py -B build-esp32c6 -D SDKCONFIG_DEFAULTS=sdkconfig.defaults flash monitor
\end{lstlisting}

\begin{lstlisting}[style=bash,caption={Building and running for the QEMU emulator (ESP32\nobreakdash-C3 target).}]
idf.py -B build-qemu -D SDKCONFIG_DEFAULTS=sdkconfig.defaults.qemu set-target esp32c3
idf.py -B build-qemu -D SDKCONFIG_DEFAULTS=sdkconfig.defaults.qemu qemu --graphics monitor
\end{lstlisting}

\section{Confirming it works}

A healthy board boot logs the firmware entering \code{app\_main}, the configured
display pins, and finally ``PRG32 runtime initialized,'' then shows the splash
\citep{prg32tutorial}. The platform also ships a smoke test that exercises the whole
workflow---environment, firmware build, cartridge build, and QEMU staging---and
prints \code{=== SMOKE TEST PASSED ===} when all is well \citep{prg32repo}.

\begin{lstlisting}[style=bash,caption={Running the end-to-end smoke test.}]
./scripts/smoke_test.sh
\end{lstlisting}

\begin{prgcheck}
Your environment is ready when, on your operating system, \code{idf.py} runs in an
ESP-IDF shell, the QEMU build opens a virtual screen reaching ``PRG32 runtime
initialized,'' and---if you have a board---the firmware flashes and shows its
splash. With that in place, every chapter and every example in this book is
buildable on your machine.
\end{prgcheck}
